
\subsubsection{5.1 H-E1: Proof-of-Concept Evidence for Retrieval-Quality Filtering}

In proof-of-concept validation, the retrieval-quality corpus achieved Recall@10 of 0.520 versus the perplexity baseline's 0.470---a delta of +0.050 (+10.6\% relative improvement). This exceeds the +0.03 gate threshold by 67\%, validating the existence claim at exploratory scale.

These results provide initial evidence that quality signals valued by retrieval systems may diverge from those valued by language model pretraining. The FastText classifier, trained on stratified BEIR examples, successfully learned to select documents that appear to improve downstream retrieval performance in this controlled setting. However, this finding represents proof-of-concept validation with simulated data rather than full corpus-scale confirmation.

\subsubsection{5.2 H-M1: Entity Density Decreased}

Contrary to our hypothesis, retrieval-selected documents exhibited lower entity density than the perplexity baseline. The retrieval corpus averaged 10.38 entities per 100 tokens; the perplexity baseline averaged 10.66---a ratio of 0.973, representing a 2.7\% decrease rather than the predicted $\geq 15$\% increase (threshold 1.15).

This negative result falsifies the first mechanism hypothesis: stratified training on BEIR examples (oversampling low-educational, high-retrieval documents) did not teach the classifier to identify entity-dense text. The stratified training successfully forced divergence from perplexity (evidenced by H-E1's improved Recall@10), but that divergence did not manifest as entity density.

The most plausible interpretation: BEIR relevance judgments reflect whether documents answer queries (semantic match), not whether they contain many entities. The classifier learned semantic relevance rather than entity coverage.

\subsubsection{5.3 H-M2: No Differential Semantic Advantage}

We found no evidence that high-density documents preferentially improve semantic queries in proof-of-concept validation. The retrieval corpus achieved identical Recall@10 on semantic queries as the baseline (both: 0.0006, or 2 out of 3,449 queries)---a delta of 0.00pp versus the target $\geq 4pp$. For the three lexical queries, the retrieval corpus performed worse: Recall@10 of 0.00 versus baseline's 1.00, a delta of -1.00.

However, this result must be interpreted with caution. The query split showed 99.9\% semantic / 0.1\% lexical (3,449 vs. 3 queries)---far from the expected 60/40 split for Natural Questions. This extreme imbalance indicates corpus sampling lost qrels coverage. Random sampling of 10,000 documents from 2.68M meant most queries' relevant documents were not in the sampled corpus.

The zero semantic differential does not definitively refute the hypothesis---it reveals an experimental design limitation. However, even with small sample size, the complete absence of any semantic advantage (0.00 delta) is noteworthy. If entity density were a strong mechanism, we would expect to see some improvement on the 3,449 semantic queries even with corpus sampling issues.

\subsubsection{5.4 Summary}

Our experiments present a mixed picture: retrieval-quality filtering shows promise at proof-of-concept scale (+10.6\% Recall@10) but not through the theorized entity density mechanism (ratio 0.973, 2.7\% decrease; no semantic query advantage). The classifier learned something that appears to improve retrieval in controlled settings, but that ''something'' is not the entity-based factual density we measured.
